% Appendix A: Technical Implementation Details
% Edge-Weighted METIS Integration

\section{Edge-Weighted METIS Integration and Implementation}
\label{appendix:implementation}

This appendix provides the technical details necessary to replicate our edge-weighted graph partitioning approach. We describe the algorithm workflow, pseudocode, METIS parameters, edge weight computation, and computational complexity.

\subsection{Algorithm Overview}

Our approach integrates demographic edge-weighting into METIS's multilevel recursive bisection algorithm. The workflow consists of four stages:

\begin{enumerate}
    \item \textbf{Graph Construction}: Convert census tract boundaries to a weighted undirected graph.
    \item \textbf{Edge Weight Assignment}: Compute edge weights reflecting both geographic boundaries and demographic composition.
    \item \textbf{Recursive Bisection}: Apply METIS's multilevel algorithm with edge weights to partition the graph into $k$ districts.
    \item \textbf{Validation}: Check contiguity, population balance, and compactness metrics.
\end{enumerate}

\subsection{Graph Construction}

\textbf{Input}: Census tract shapefiles with demographics (total population, minority population by race/ethnicity).

\textbf{Output}: Undirected graph $G = (V, E, w_V, w_E)$ where:
\begin{itemize}
    \item $V$: Set of vertices (census tracts)
    \item $E$: Set of edges (shared boundaries between adjacent tracts)
    \item $w_V: V \to \mathbb{R}^+$: Vertex weights (tract populations)
    \item $w_E: E \to \mathbb{R}^+$: Edge weights (computed by edge-weighting function)
\end{itemize}

\textbf{Adjacency Detection}: Two tracts $u, v \in V$ are adjacent (share edge $(u,v) \in E$) if their boundaries share a non-zero-length segment. We use GeoPandas' spatial operations:

\begin{verbatim}
def build_adjacency_graph(tracts_gdf):
    """
    Build adjacency graph from census tract geometries.

    Parameters:
        tracts_gdf: GeoDataFrame with columns ['GEOID', 'geometry',
                    'total_pop', 'minority_pop']

    Returns:
        G: NetworkX graph with vertex/edge weights
    """
    import networkx as nx
    from shapely.ops import shared_paths

    G = nx.Graph()

    # Add vertices with population weights
    for idx, tract in tracts_gdf.iterrows():
        G.add_node(
            tract['GEOID'],
            weight=tract['total_pop'],
            minority_pct=tract['minority_pop'] / tract['total_pop']
        )

    # Detect adjacency via spatial intersection
    for i, tract_i in tracts_gdf.iterrows():
        for j, tract_j in tracts_gdf.iterrows():
            if i < j:  # Avoid duplicate edges
                if tract_i.geometry.touches(tract_j.geometry):
                    G.add_edge(tract_i['GEOID'], tract_j['GEOID'])

    return G
\end{verbatim}

\textbf{Note}: For large states (1000+ tracts), we use spatial indexing (R-tree) to accelerate adjacency detection from $O(n^2)$ to $O(n \log n)$.

\subsection{Edge Weight Computation}

Edge weights encode both geographic proximity and demographic composition. The weight $w(u,v)$ for edge $(u,v) \in E$ is computed as:

\[
w(u,v) = \ell(u,v) \times \alpha(u,v)
\]

where:
\begin{itemize}
    \item $\ell(u,v)$: Geographic length of shared boundary between tracts $u$ and $v$ (in meters).
    \item $\alpha(u,v)$: Demographic scaling factor reflecting minority composition.
\end{itemize}

\subsubsection{Demographic Scaling Factor $\alpha(u,v)$}

The scaling factor $\alpha(u,v)$ amplifies edge weights for minority-minority edges:

\[
\alpha(u,v) =
\begin{cases}
\beta & \text{if } \min(p_u, p_v) \geq \tau \\
1.0 & \text{otherwise}
\end{cases}
\]

where:
\begin{itemize}
    \item $p_u, p_v$: Minority percentages in tracts $u$ and $v$
    \item $\tau$: Minority threshold (e.g., 40\%, 45\%, 50\%)
    \item $\beta$: Edge weight scaling factor (e.g., 1×, 5×, 10×)
\end{itemize}

\textbf{Interpretation}: If both tracts $u$ and $v$ have minority percentages $\geq \tau$, the edge weight is multiplied by $\beta$, making it more ``expensive'' for METIS to cut. This nudges the algorithm toward keeping minority clusters together without hard constraints.

\textbf{Example}: Alabama with $\beta = 5$, $\tau = 45\%$:
\begin{itemize}
    \item Edge between two tracts with 55\% Black population each: $w = \ell \times 5.0$
    \item Edge between 55\% Black tract and 30\% Black tract: $w = \ell \times 1.0$
    \item Edge between two 30\% Black tracts: $w = \ell \times 1.0$
\end{itemize}

\subsubsection{Pseudocode for Edge Weighting}

\begin{verbatim}
def compute_edge_weights(G, beta, tau):
    """
    Assign edge weights based on demographic composition.

    Parameters:
        G: NetworkX graph with 'minority_pct' node attributes
        beta: Scaling factor for minority-minority edges
        tau: Minority threshold (e.g., 0.45 for 45%)

    Returns:
        G: Graph with updated edge weights
    """
    for u, v in G.edges():
        # Get minority percentages
        p_u = G.nodes[u]['minority_pct']
        p_v = G.nodes[v]['minority_pct']

        # Compute boundary length (meters)
        length = compute_boundary_length(u, v)

        # Apply demographic scaling
        if min(p_u, p_v) >= tau:
            weight = length * beta
        else:
            weight = length * 1.0

        G[u][v]['weight'] = weight

    return G
\end{verbatim}

\subsection{METIS Multilevel Recursive Bisection}

We use METIS 5.1.0's multilevel recursive bisection algorithm (\texttt{METIS\_PartGraphRecursive}) to partition the graph into $k$ districts.

\subsubsection{METIS Parameters}

The following parameters control METIS's behavior:

\begin{itemize}
    \item \textbf{objtype}: \texttt{METIS\_OBJTYPE\_CUT} (minimize edge cut, not communication volume)
    \item \textbf{ctype}: \texttt{METIS\_CTYPE\_SHEM} (sorted heavy-edge matching for coarsening)
    \item \textbf{iptype}: \texttt{METIS\_IPTYPE\_GROW} (greedy graph growing for initial partitioning)
    \item \textbf{rtype}: \texttt{METIS\_RTYPE\_GREEDY} (greedy boundary refinement)
    \item \textbf{ufactor}: 5 (population imbalance tolerance: 0.5\%)
    \item \textbf{niter}: 10 (number of refinement iterations)
    \item \textbf{seed}: 42 (random seed for reproducibility)
    \item \textbf{dbglvl}: 0 (no debug output)
\end{itemize}

\subsubsection{Population Balance Constraint}

METIS enforces population balance through vertex weights. Each tract's population becomes its vertex weight $w_V(v)$, and METIS ensures each district has total population within $(1 \pm \epsilon)$ of the target, where $\epsilon = 0.005$ (0.5\%).

\textbf{Target Population}:
\[
P_{\text{target}} = \frac{\sum_{v \in V} w_V(v)}{k}
\]

\textbf{Balance Constraint}:
\[
0.995 \times P_{\text{target}} \leq P_i \leq 1.005 \times P_{\text{target}} \quad \forall i \in \{1, \ldots, k\}
\]

where $P_i$ is the total population of district $i$.

\subsubsection{Integration with Python}

We use PyMETIS, a Python wrapper for METIS, with the following interface:

\begin{verbatim}
import pymetis
import networkx as nx

def partition_with_metis(G, k, beta, tau):
    """
    Partition graph into k districts using edge-weighted METIS.

    Parameters:
        G: NetworkX graph with vertex/edge weights
        k: Number of districts
        beta: Edge weight scaling factor
        tau: Minority threshold

    Returns:
        partition: Dict mapping tract GEOID to district ID
    """
    # Apply edge weighting
    G = compute_edge_weights(G, beta, tau)

    # Convert to METIS format
    adjacency = [list(G.neighbors(node)) for node in G.nodes()]
    node_weights = [G.nodes[node]['weight'] for node in G.nodes()]
    edge_weights = []
    for u in G.nodes():
        edge_weights.append([G[u][v]['weight'] for v in G.neighbors(u)])

    # Call METIS
    edgecut, parts = pymetis.part_graph(
        k,  # Number of parts
        adjacency=adjacency,
        vweights=node_weights,
        eweights=edge_weights,
        recursive=True,  # Use recursive bisection
        seed=42
    )

    # Map back to GEOIDs
    partition = {node: parts[i] for i, node in enumerate(G.nodes())}

    return partition, edgecut
\end{verbatim}

\subsection{Recursive Bisection Workflow}

METIS's recursive bisection algorithm partitions a graph with $k$ districts through $\lceil \log_2 k \rceil$ recursive bisection steps:

\begin{enumerate}
    \item \textbf{Coarsening}: Repeatedly match and merge vertices to create a hierarchy of smaller graphs $G_0 \to G_1 \to \cdots \to G_L$.
    \item \textbf{Initial Partitioning}: Partition the coarsest graph $G_L$ into 2 parts using greedy graph growing.
    \item \textbf{Uncoarsening}: Project the 2-way partition back through the hierarchy $G_L \to G_{L-1} \to \cdots \to G_0$, refining boundaries at each level using Kernighan-Lin / Fiduccia-Mattheyses local search.
    \item \textbf{Recursive Application}: Apply steps 1-3 to each partition recursively until $k$ parts are obtained.
\end{enumerate}

\textbf{Edge Weight Integration}: Edge weights $w_E$ are preserved throughout coarsening. When vertices $u$ and $v$ are merged into supervertex $w$, edge weights are summed:
\[
w_E(w, x) = w_E(u, x) + w_E(v, x) \quad \forall x \in \text{neighbors}(w)
\]

This ensures that demographic considerations (encoded in edge weights) influence partitioning at all levels of the hierarchy.

\subsection{Computational Complexity}

\subsubsection{Time Complexity}

\begin{itemize}
    \item \textbf{Graph Construction}: $O(n \log n)$ using R-tree spatial indexing, where $n = |V|$ (number of tracts).
    \item \textbf{Edge Weighting}: $O(m)$, where $m = |E|$ (number of edges). For planar graphs (tract adjacency), $m \approx 3n$.
    \item \textbf{METIS Partitioning}: $O((n + m) \log k)$ for recursive bisection with $k$ districts and multilevel refinement.
    \item \textbf{Total}: $O(n \log n + m \log k) = O(n \log n)$ since $m = O(n)$ and $k \ll n$ (typically $k \leq 14$ for congressional redistricting).
\end{itemize}

\subsubsection{Space Complexity}

\begin{itemize}
    \item \textbf{Graph Storage}: $O(n + m) = O(n)$ for adjacency lists.
    \item \textbf{METIS Workspace}: $O(n)$ for coarsening/uncoarsening buffers.
    \item \textbf{Total}: $O(n)$ linear in the number of tracts.
\end{itemize}

\subsubsection{Empirical Runtimes}

\begin{table}[H]
\centering
\begin{tabular}{lcccc}
\toprule
State & Tracts & Districts & Runtime (s) & Runtime/Tract (ms) \\
\midrule
Vermont & 181 & 1 & 2.3 & 12.7 \\
Delaware & 218 & 1 & 2.8 & 12.8 \\
Alabama & 1186 & 7 & 14.2 & 12.0 \\
Georgia & 1973 & 14 & 28.5 & 14.4 \\
South Carolina & 1106 & 7 & 13.8 & 12.5 \\
Louisiana & 1212 & 6 & 15.1 & 12.5 \\
Mississippi & 836 & 4 & 9.7 & 11.6 \\
California & 9036 & 52 & 132.4 & 14.7 \\
Texas & 5823 & 38 & 88.9 & 15.3 \\
\midrule
\textbf{Mean} & --- & --- & --- & \textbf{13.2 ms/tract} \\
\bottomrule
\end{tabular}
\caption{Empirical runtimes for edge-weighted METIS partitioning on a single core (Intel Core i7-10700K @ 3.8 GHz, 32 GB RAM). Runtime scales approximately linearly with tract count, averaging 13.2 ms per tract.}
\label{tab:runtimes}
\end{table}

\subsection{Software and Reproducibility}

\subsubsection{Dependencies}

Our implementation requires:
\begin{itemize}
    \item Python 3.9+
    \item METIS 5.1.0 (compiled with 32-bit vertex indices)
    \item PyMETIS 2022.1
    \item NetworkX 2.8+
    \item GeoPandas 0.11+
    \item Shapely 1.8+
    \item Pandas 1.4+
    \item NumPy 1.22+
\end{itemize}

\subsubsection{Installation}

METIS must be compiled from source with specific flags to support large graphs:

\begin{verbatim}
# Download METIS 5.1.0
wget http://glaros.dtc.umn.edu/gkhome/fetch/sw/metis/metis-5.1.0.tar.gz
tar -xzf metis-5.1.0.tar.gz
cd metis-5.1.0

# Compile with 32-bit indices (supports up to 2^32 vertices)
make config shared=1 cc=gcc prefix=/usr/local
make install

# Install PyMETIS
pip install pymetis
\end{verbatim}

\subsubsection{Code Repository}

Full implementation code, data processing scripts, and replication instructions are available at:

\begin{center}
\texttt{[Repository URL to be provided upon publication]}
\end{center}

The repository includes:
\begin{itemize}
    \item Graph construction pipeline (adjacency detection, edge weighting)
    \item METIS wrapper with parameter grid search
    \item Compactness and VRA metric computation
    \item Visualization generation (Pareto frontiers, district maps)
    \item Example notebooks for replicating key results
\end{itemize}

\subsubsection{Data Availability}

Census tract shapefiles and demographics are publicly available:
\begin{itemize}
    \item \textbf{Geometries}: U.S. Census Bureau TIGER/Line Shapefiles (2020)
    \begin{itemize}
        \item URL: \url{https://www.census.gov/geographies/mapping-files/time-series/geo/tiger-line-file.2020.html}
    \end{itemize}
    \item \textbf{Demographics}: Census 2020 Redistricting Data (Public Law 94-171)
    \begin{itemize}
        \item URL: \url{https://www.census.gov/programs-surveys/decennial-census/about/rdo/summary-files.html}
    \item Variables: Total population (P1), Race by ethnicity (P2, P3)
    \end{itemize}
\end{itemize}

Processed datasets (tract-level minority percentages, adjacency graphs) are included in the code repository to facilitate replication without reprocessing raw census files.

\subsection{Parameter Sensitivity Analysis}

We tested 105 configurations across 5 states, varying two parameters:

\begin{itemize}
    \item \textbf{Edge weight scaling $\beta$}: \{1×, 5×, 10×, 50×, 100×, 1000×\}
    \item \textbf{Minority threshold $\tau$}: \{30\%, 35\%, 40\%, 45\%, 50\%, 55\%\}
\end{itemize}

\textbf{Key Findings}:
\begin{itemize}
    \item \textbf{Alabama}: Optimal at $\beta = 5$, $\tau = 45\%$ (2 MM, +3.2\% PP, $-9.3$\% edge cut). Higher $\beta$ (10×, 50×) achieves same MM count with slightly worse compactness.
    \item \textbf{Georgia}: Plateau at $\beta = 5$, $\tau = 50\%-55\%$ (6 MM, +22\% PP). Further increasing $\beta$ yields no additional MM districts.
    \item \textbf{South Carolina}: No configuration achieves $\geq 2$ MM districts, even with $\beta = 1000$, $\tau = 30\%$, confirming infeasibility.
    \item \textbf{Threshold Sensitivity}: Lower $\tau$ (30\%-40\%) increases MM district counts but degrades compactness. $\tau = 45\%-50\%$ balances VRA compliance with compactness.
    \item \textbf{Scaling Sensitivity}: Moderate $\beta$ (5×-10×) suffices for most states. Extreme scaling (100×-1000×) provides no additional benefit and risks numerical instability.
\end{itemize}

\textbf{Recommendation}: Use $\beta = 5$, $\tau = 45\%$ as default parameters for VRA-compliant redistricting, with parameter sweeps to explore Pareto frontiers.

\subsection{Validation and Robustness Checks}

\subsubsection{Contiguity Validation}

After METIS partitioning, we verify that all districts are contiguous:

\begin{verbatim}
def validate_contiguity(partition, G):
    """Check that each district forms a connected component."""
    for district_id in set(partition.values()):
        tracts = [t for t, d in partition.items() if d == district_id]
        subgraph = G.subgraph(tracts)
        if not nx.is_connected(subgraph):
            raise ValueError(f"District {district_id} is non-contiguous")
    return True
\end{verbatim}

All 105 configurations across 5 states produced contiguous districts, confirming that edge-weighting does not compromise contiguity.

\subsubsection{Population Balance Validation}

We verify population deviation $\leq 0.5\%$ for all districts:

\[
\text{Deviation}_i = \left| \frac{P_i - P_{\text{target}}}{P_{\text{target}}} \right| \times 100\%
\]

All configurations satisfied deviation $\leq 0.5\%$, with median deviation $< 0.1\%$.

\subsubsection{Seed Sensitivity}

We tested 10 random seeds for Alabama's optimal configuration (5×@45\%):
\begin{itemize}
    \item \textbf{MM districts}: 2/7 in all 10 runs (no variation)
    \item \textbf{Edge cut}: $254 \pm 2$ (coefficient of variation: 0.8\%)
    \item \textbf{Polsby-Popper}: $0.242 \pm 0.003$ (CV: 1.2\%)
\end{itemize}

Results are robust to random seed variation, indicating stable Pareto frontiers.

\subsection{Comparison to Multi-Constraint Approach}

Our edge-weighted approach differs fundamentally from METIS's multi-constraint partitioning (\texttt{METIS\_PartGraphKway} with multiple vertex weight vectors):

\begin{itemize}
    \item \textbf{Multi-Constraint}: Each vertex has multiple weights (e.g., total population, minority population). METIS enforces balance on \emph{all} weight dimensions simultaneously. For redistricting, this attempts to balance both total population and minority population across districts, forcing similar minority percentages in all districts.
    \item \textbf{Edge-Weighted}: Each vertex has a single weight (total population). Demographic considerations enter through \emph{edge weights}, not vertex constraints. This allows flexible clustering---districts may have very different minority percentages, respecting natural clustering.
\end{itemize}

\textbf{Alabama Comparison (Paper 4 Result)}:
\begin{itemize}
    \item \textbf{Multi-Constraint}: 0 MM districts, edge cut 280, Polsby-Popper 0.234
    \item \textbf{Edge-Weighted (5×@45\%)}: 2 MM districts, edge cut 254, Polsby-Popper 0.242
    \item \textbf{Conclusion}: Edge-weighted approach dominates on both VRA compliance and compactness.
\end{itemize}

\textbf{Mechanistic Explanation}: Multi-constraint forces all districts toward similar minority percentages (~36-37\% in Alabama), preventing formation of $\geq 50\%$ MM districts. Edge-weighting allows heterogeneity---some districts reach 51-52\% minority (MM), others drop to 25-30\% (non-MM), better reflecting natural clustering.
